% Chapter 1 - Introduction
\section{Introduction}
\label{chap:introduction}

\subsection{Motivation}
\label{sec:motivation}

For decades, the security of widely deployed public-key cryptosystems such as RSA
and Elliptic Curve Cryptography (ECC) has rested on the assumed intractability of
integer factorization and the discrete logarithm problem for classical computers.
\citet{shor1997} showed that both problems can be solved in polynomial time on a
sufficiently large quantum computer, meaning that a cryptographically relevant
quantum computer would render these widely used schemes insecure. Although such a
machine does not yet exist, the risk that data encrypted today could be decrypted
retroactively once quantum computers become available -- commonly referred to as
``store now, decrypt later'' -- has already prompted standardization bodies to act.
In August 2024, the National Institute of Standards and Technology (NIST) published
the first three finalized post-quantum cryptography (PQC) standards
\citep{nist_pqc}, and the German Federal Office for Information Security (BSI)
recommends that organizations begin migrating to quantum-safe algorithms without
delay \citep{bsi_broschuere}. Post-quantum cryptography has thus moved from a
theoretical research topic to an active engineering concern.

Among the algorithm families studied for this purpose, code-based cryptography is
one of the oldest. The McEliece cryptosystem, introduced by \citet{mceliece1978}
almost five decades ago, has withstood cryptanalysis ever since and is still
regarded as one of the most conservative choices for post-quantum encryption
\citep{classic_mceliece_slides2024}. Its most recent NIST submission, Classic
McEliece, was not selected for standardization in the fourth round of the NIST PQC
process -- not because of a security weakness, but primarily because of its
comparatively large public keys and to avoid diverging from a parallel ISO
standardization effort \citep{nist_ir8545}. This makes McEliece a notable case: a
cryptosystem that is simultaneously among the most trusted and among the least
understood by students, since its security relies on coding-theoretic concepts
(linear codes, Goppa codes) that are typically not part of an undergraduate
curriculum \citep{overbeck2009codebased}.

Interactive, browser-based learning tools have proven effective at making other
areas of cryptography more approachable. CrypTool Online (CTO) already provides
such visualizations for a range of classical and modern algorithms. However, no
comparable visualization currently exists for the McEliece cryptosystem within CTO,
leaving a gap between the system's practical and historical importance and the
resources available to teach it.

\subsection{Research Question}
\label{sec:researchQuestion}

Building on this motivation, this thesis addresses the following research question:

\begin{quote}
	\itshape How can the McEliece cryptosystem be implemented and visualized within
	CrypTool Online to improve the understanding of code-based post-quantum
	cryptography for students?
\end{quote}

\subsection{Structure of the Thesis}
\label{sec:thesisStructure}

The remainder of this thesis is structured as follows. \autoref{chap:cryptoFundamentals}
briefly introduces the cryptographic fundamentals needed as a common frame of
reference. \autoref{chap:qcAndPqc} covers the fundamentals of quantum computing and
surveys the post-quantum cryptography landscape, positioning McEliece within it.
\autoref{chap:mceliece} presents the McEliece cryptosystem itself, from the
underlying coding theory to key generation, encryption, and decryption.
\autoref{chap:ctoArchitecture} introduces CrypTool Online and its plugin
architecture as the technical platform for this work, and reviews related
educational tools. \autoref{chap:educationalConcept} develops the educational and
didactic concept underlying the planned visualization.
\autoref{chap:requirements} derives the requirements for the plugin from the
problem statement, objectives, and scope of this thesis. \autoref{chap:pluginDesign}
presents the design of the McEliece learning plugin, and
\autoref{chap:implementation} describes its implementation within CTO.
\autoref{chap:evaluation} evaluates the resulting artifact, and
\autoref{chap:conclusion} concludes the thesis and outlines future work.
